\section{Planets Queries II}
\label{sec:cses-planets-queries-ii}

\problemstatement{Each of $n$ planets has one teleporter. Answer $q$
queries: the minimum number of teleports to go from planet $a$ to planet
$b$, or $-1$ if impossible.}

\begin{patternbox}
A functional graph is ``rho-shaped'': trees of tail nodes feeding into
cycles. Decompose it into cycles and the depth of each node to its cycle,
then a query splits into a \textbf{tail case} (answered by binary lifting)
and a \textbf{cycle case} (answered by modular arithmetic on the cycle).
\end{patternbox}

\subsection*{Tails into cycles}

Peel nodes of in-degree $0$ repeatedly; whatever remains lies on cycles.
Label each cycle, giving its nodes positions $0,1,\ldots$ and recording the
cycle length. Each tail node gets a $\textit{depth}$ (steps to reach its
cycle) and an $\textit{entry}$ (the cycle node it first lands on). For a
query $(a,b)$:
\begin{itemize}
  \item \textbf{Tail case:} if $b$ lies on the path from $a$ into the
        cycle, then $\textit{depth}[a]\ge\textit{depth}[b]$ and jumping
        $\textit{depth}[a]-\textit{depth}[b]$ steps (binary lifting) from
        $a$ lands on $b$; that jump count is the answer.
  \item \textbf{Cycle case:} if $b$ is on the same cycle $a$ reaches, walk
        $\textit{depth}[a]$ steps to $a$'s entry, then around the cycle:
        $\textit{depth}[a]+((\textit{pos}[b]-\textit{pos}[\textit{entry}[a]])
        \bmod L)$.
\end{itemize}
Take whichever applies; otherwise $-1$.

\begin{ideabox}[Split the Query by Where b Lives]
The rho structure means $b$ is reachable from $a$ only if it is on $a$'s
tail before the cycle, or on the cycle $a$ enters. Binary lifting verifies
the tail case exactly; modular position arithmetic handles the cycle case.
\end{ideabox}

\subsection*{Algorithm}

\begin{enumerate}
  \item Build binary lifting for the successor function.
  \item Peel to find cycle nodes; assign cycle ids, positions, lengths;
        compute $\textit{depth}$ and $\textit{entry}$ for tail nodes.
  \item Per query, evaluate the tail and cycle cases; output the minimum or
        $-1$.
\end{enumerate}

\subsection*{Dry run}

\dryrun{$\textit{succ}=[2,3,1,3]$: $1\!\to\!2\!\to\!3\!\to\!1$ cycle, $4\!
\to\!3$ tail.}

\begin{itemize}
  \item Query $(4,1)$: $4$ enters the cycle at $3$ in $1$ step, then $3\to
        1$ is $2$ more steps around the cycle. Answer $3$.
  \item Query $(1,4)$: $4$ is a tail node not reachable from the cycle --
        $-1$.
\end{itemize}

\subsection*{C++ implementation}

\begin{lstlisting}
#include <bits/stdc++.h>
using namespace std;
const int LOG = 20;

int main() {
    int n, q;
    cin >> n >> q;
    vector<int> succ(n + 1);
    vector<vector<int>> up(LOG, vector<int>(n + 1));
    vector<int> indeg(n + 1, 0);
    for (int i = 1; i <= n; i++) { cin >> succ[i]; up[0][i] = succ[i]; indeg[succ[i]]++; }
    for (int j = 1; j < LOG; j++)
        for (int i = 1; i <= n; i++) up[j][i] = up[j - 1][up[j - 1][i]];

    auto jump = [&](int x, int k) {
        for (int j = 0; j < LOG; j++) if (k & (1 << j)) x = up[j][x];
        return x;
    };

    // peel tree nodes; remaining are cycle nodes
    vector<int> deg = indeg, order;
    queue<int> pq;
    for (int i = 1; i <= n; i++) if (deg[i] == 0) pq.push(i);
    vector<bool> removed(n + 1, false);
    while (!pq.empty()) {
        int u = pq.front(); pq.pop();
        removed[u] = true; order.push_back(u);
        if (--deg[succ[u]] == 0) pq.push(succ[u]);
    }

    vector<int> depth(n + 1, 0), entry(n + 1, 0), lead(n + 1, 0), pos(n + 1, 0);
    vector<bool> isCyc(n + 1, false);
    vector<int> cycLen(n + 1, 0);
    int cid = 0;
    for (int i = 1; i <= n; i++) {
        if (removed[i] || isCyc[i]) continue;
        cid++; int c = i, len = 0;
        for (int x = i; !isCyc[x]; x = succ[x]) {
            isCyc[x] = true; pos[x] = len++; lead[x] = cid;
            entry[x] = x; depth[x] = 0;
        }
        cycLen[cid] = len;
    }
    for (int idx = (int)order.size() - 1; idx >= 0; idx--) {   // reverse peel order
        int u = order[idx];
        depth[u] = depth[succ[u]] + 1;
        entry[u] = entry[succ[u]];
        lead[u] = lead[succ[u]];
    }

    while (q--) {
        int a, b;
        cin >> a >> b;
        long long ans = -1;
        if (a == b) { cout << 0 << "\n"; continue; }
        // tail case: b on a's path into the cycle
        if (depth[a] >= depth[b]) {
            int steps = depth[a] - depth[b];
            if (jump(a, steps) == b) ans = steps;
        }
        // cycle case: b on the same cycle a reaches
        if (isCyc[b] && lead[b] == lead[a]) {
            int L = cycLen[lead[a]];
            long long steps = depth[a] + ((pos[b] - pos[entry[a]]) % L + L) % L;
            if (ans == -1 || steps < ans) ans = steps;
        }
        cout << ans << "\n";
    }
    return 0;
}
\end{lstlisting}

\begin{complexitybox}
\textbf{Time Complexity.} $O(n\log n)$ preprocessing, $O(\log n)$ per query.
\textbf{Space Complexity.} $O(n\log n)$ for the lifting table.
\end{complexitybox}

\begin{takeawaybox}
Distances in a functional graph split by structure: binary lifting settles
the tree-tail case, modular cycle arithmetic settles the cycle case. Seeing
the rho shape -- tails feeding cycles -- is the key to organising the query
logic.
\end{takeawaybox}
